% Fuente bilingüe común: las fórmulas y pruebas estructurales son idénticas.
\clearpage
\section{\HMTLang{Transport symplectique de l'ellipse d'action}{Transport symplectique de l'ellipse d'action}}
\label{delta22:ii:transporte}
\HMTLang{L'ellipse précédemment construite permet de distinguer la mesure aréale de l'action, l'anisotropie angulaire et le taux d'évolution. Cette distinction détermine un transport exact entre réalisations elliptiques et fixe le terme de connexion de son expression hamiltonienne. L'antécédent est la paire angulaire exprimée par APP--TRIT--TPK, après la génération régionale et la sélection conjointe de \(K\) ; les opérations suivantes agissent sur ces sorties et conservent leurs étiquettes d'orientation.}{L'ellipse construite ci-dessus distingue la mesure aréale de l'action, l'anisotropie angulaire et le taux d'évolution. Cette distinction détermine un transport exact entre réalisations elliptiques et fixe le terme de connexion de son expression hamiltonienne. Son antécédent est la paire angulaire publiée par APP--TRIT--TPK après la génération régionale et la sélection conjointe de \(K\) ; les opérations suivantes agissent sur ces sorties et conservent leurs étiquettes d'orientation.}

\subsection{\HMTLang{Forme quadratique et réalisation inductive--capacitive}{Forme quadratique et réalisation inductive--capacitive}}
\HMTLang{Nous écrivons \(\zeta=\eta_{\rm el}=\operatorname{artanh}(C^*/A)\), avec les deux coordonnées angulaires dans la même unité et \(|C^*/A|<1\). Le symbole \(\eta_{\rm ret}\) est réservé au coefficient d'action. Pour la section positive \(\hbar=\eta_{\rm ret}s_0\), \(s_0=10^{-34}\mathcal U_S\), les coordonnées équilibrées de dimension racine carrée d'action sont}{Écrivons \(\zeta=\eta_{\rm el}=\operatorname{artanh}(C^*/A)\), avec les deux coordonnées angulaires dans la même unité et \(|C^*/A|<1\). Le symbole \(\eta_{\rm ret}\) est réservé au coefficient d'action. Pour la section positive \(\hbar=\eta_{\rm ret}s_0\), \(s_0=10^{-34}\mathcal U_S\), les coordonnées équilibrées, de dimension racine carrée d'action, sont}
\[
 Q=\sqrt{2\hbar}\,e^{\zeta/2}\cos\theta,
 \qquad P=\sqrt{2\hbar}\,e^{-\zeta/2}\sin\theta.
\]
\HMTLang{Le produit des demi-axes est \(2\hbar\), l'aire totale est \(2\pi\hbar=h\) et l'aire paramétrique cumulée est \(\hbar\theta\). La phase \(\theta\) est le paramètre de cette représentation elliptique. En conservant le transducteur \(Z_*>0\) et la fréquence \(\omega>0\) de la réalisation précédente, on a}{Le produit des demi-axes est \(2\hbar\), l'aire totale est \(2\pi\hbar=h\) et l'aire paramétrique cumulée est \(\hbar\theta\). La phase \(\theta\) est le paramètre de cette représentation elliptique. En conservant le transducteur \(Z_*>0\) et la fréquence \(\omega>0\) de la réalisation précédente, on obtient}
\[
 L_\zeta=\frac{Z_*}{\omega}e^{-\zeta},\qquad
 C_\zeta=\frac{e^\zeta}{Z_*\omega},\qquad
 Q=\sqrt{Z_*}\,q,\quad P=\frac{\Phi}{\sqrt{Z_*}}.
\]
\begin{proposition}[\HMTLang{Composition énergétique de l'ellipse}{Composition énergétique de l'ellipse}]
\[
 H_0=\frac{\omega}{2}(e^{-\zeta}Q^2+e^\zeta P^2),\quad
 E_{\rm el}=\hbar\omega\cos^2\theta,\quad
 E_{\rm mag}=\hbar\omega\sin^2\theta.
\]
\HMTLang{Sur le contour, \(H_0=E_{\rm el}+E_{\rm mag}=\hbar\omega\).}{Sur la frontière, \(H_0=E_{\rm el}+E_{\rm mag}=\hbar\omega\).}
\end{proposition}
\begin{proof}
\HMTLang{La substitution de \(q=Q/\sqrt{Z_*}\), \(\Phi=\sqrt{Z_*}P\) dans \(q^2/(2C_\zeta)+\Phi^2/(2L_\zeta)\) produit \(H_0\). La substitution ultérieure de la paramétrisation elliptique donne les deux contributions et leur somme. L'action et la forme précèdent le circuit qui les réalise.}{La substitution de \(q=Q/\sqrt{Z_*}\), \(\Phi=\sqrt{Z_*}P\) dans \(q^2/(2C_\zeta)+\Phi^2/(2L_\zeta)\) donne \(H_0\). La substitution ultérieure de la paramétrisation elliptique donne les deux contributions et leur somme. L'action et la forme précèdent le circuit qui les réalise.}
\end{proof}
\HMTLang{L'impédance de cette réalisation satisfait \(\sqrt{L_\zeta/C_\zeta}/Z_*=e^{-\zeta}\). L'impédance constitutive du vide utilise le quotient de ses lecteurs \(90/120\) ; chaque impédance conserve sa définition et sa transduction.}{L'impédance de cette réalisation satisfait \(\sqrt{L_\zeta/C_\zeta}/Z_*=e^{-\zeta}\). L'impédance constitutive du vide utilise le rapport de ses lecteurs \(90/120\) ; chaque impédance conserve sa définition et sa transduction.}

\subsection{\HMTLang{Invariant de transport et connexion hamiltonienne}{Invariant de transport et connexion hamiltonienne}}
\HMTLang{Soient}{Let}
\[
 S_\zeta=\operatorname{diag}(e^{\zeta/2},e^{-\zeta/2}),\quad
 g_\zeta=S_\zeta^{-\mathsf T}S_\zeta^{-1},\quad
 J_0=\begin{pmatrix}0&-1\\1&0\end{pmatrix},\quad
 R(-\theta)=e^{-\theta J_0}.
\]
\begin{theorem}[\HMTLang{Conservation exacte de l'action quadratique}{Conservation exacte de l'action quadratique}]
\HMTLang{Pour une suite de formes \(\zeta_n\) et de phases \(\theta_n\), le transport}{Pour une suite de formes \(\zeta_n\) et de phases \(\theta_n\), le transport}
\[
 X_{n+1}=T_nX_n,\qquad
 T_n=S_{\zeta_{n+1}}R(-\theta_n)S_{\zeta_n}^{-1}
\]
\HMTLang{a pour déterminant un et conserve}{a pour déterminant un et conserve}
\[
 I_n=\tfrac12X_n^{\mathsf T}g_{\zeta_n}X_n,
 \qquad T_n^{\mathsf T}g_{\zeta_{n+1}}T_n=g_{\zeta_n}.
\]
\end{theorem}
\begin{proof}
\HMTLang{Les déterminants de \(S_\zeta\) et \(R\) valent un. Dans le produit quadratique, les facteurs \(S_{\zeta_{n+1}}\) se simplifient avec leurs inverses ; \(R^{\mathsf T}R=I\) laisse \(S_{\zeta_n}^{-\mathsf T}S_{\zeta_n}^{-1}\). Cette identité prouve la conservation à chaque pas et, par composition, pour toute longueur. Une suite exprimée par une route TPK est transportée avec cette route et ses étiquettes.}{Les déterminants de \(S_\zeta\) et \(R\) valent un. Dans le produit quadratique, les facteurs \(S_{\zeta_{n+1}}\) se simplifient avec leurs inverses ; \(R^{\mathsf T}R=I\) laisse \(S_{\zeta_n}^{-\mathsf T}S_{\zeta_n}^{-1}\). Cette identité prouve la conservation à chaque pas et, par composition, pour toute longueur. Une suite publiée par une route TPK est transportée avec cette route et ses étiquettes.}
\end{proof}
\begin{theorem}[\HMTLang{Générateur différentiable et terme mixte}{Générateur différentiable et terme mixte}]
\HMTLang{Si \(\zeta\) est de classe \(C^1\), \(\omega>0\) est continue et \(X=S_\zeta Y\), \(\dot Y=-\omega J_0Y\), alors}{Si \(\zeta\) est de classe \(C^1\), \(\omega>0\) est continue et \(X=S_\zeta Y\), \(\dot Y=-\omega J_0Y\), alors}
\[
 \dot Q=\tfrac12\dot\zeta Q+\omega e^\zeta P,
 \qquad \dot P=-\omega e^{-\zeta}Q-\tfrac12\dot\zeta P.
\]
\HMTLang{Avec la convention \(\dot Q=\partial_PH\), \(\dot P=-\partial_QH\), son Hamiltonien, à un terme additif dépendant du temps près, est}{Avec la convention \(\dot Q=\partial_PH\), \(\dot P=-\partial_QH\), son Hamiltonien, à une fonction additive du temps près, est}
\begin{equation}
 H_{\rm tot}=\frac\omega2(e^{-\zeta}Q^2+e^\zeta P^2)
             +\frac{\dot\zeta}{2}QP.
 \label{delta22:ii:hamiltoniano}
\end{equation}
\HMTLang{L'invariant \(I=(e^{-\zeta}Q^2+e^\zeta P^2)/2\) est conservé exactement.}{L'invariant \(I=(e^{-\zeta}Q^2+e^\zeta P^2)/2\) est conservé exactement.}
\end{theorem}
\begin{proof}
\HMTLang{La dérivation de \(X=S_\zeta Y\) produit \(\dot S_\zeta S_\zeta^{-1}=\operatorname{diag}(\dot\zeta/2,-\dot\zeta/2)\) et le terme rotationnel conjugué. Les dérivées partielles de \eqref{delta22:ii:hamiltoniano} reproduisent les deux équations. Dans \(\dot I\), les contributions explicites \(-\dot\zeta e^{-\zeta}Q^2/2+\dot\zeta e^\zeta P^2/2\) annulent celles provenant de la diagonale du générateur ; les termes croisés rotationnels s'annulent mutuellement.}{La dérivation de \(X=S_\zeta Y\) donne \(\dot S_\zeta S_\zeta^{-1}=\operatorname{diag}(\dot\zeta/2,-\dot\zeta/2)\) et le terme rotationnel conjugué. Les dérivées partielles de \eqref{delta22:ii:hamiltoniano} reproduisent les deux équations. Dans \(\dot I\), les contributions explicites \(-\dot\zeta e^{-\zeta}Q^2/2+\dot\zeta e^\zeta P^2/2\) annulent celles de la diagonale du générateur ; les termes croisés rotationnels s'annulent mutuellement.}
\end{proof}
\HMTLang{Dans la paramétrisation elliptique, \(QP=I\sin2\theta\) et \(\dot\theta=-\omega\). Par conséquent, \(H_{\rm tot}=I\omega+(I\dot\zeta/2)\sin2\theta\). La matrice quadratique a pour déterminant \(\omega^2-\dot\zeta^2/4\) : elle est définie positive exactement pour \(|\dot\zeta|<2\omega\), et semi-définie positive dans le cas d'égalité. La conservation de \(I\) et la dépendance temporelle de \(H_{\rm tot}\) expriment des propriétés distinctes. À titre de contrôle algébrique, omettre le terme mixte produit \(\dot I=\dot\zeta(e^\zeta P^2-e^{-\zeta}Q^2)/2\). Dans la représentation quantique correspondante, l'expression symétrique du terme mixte est \(\dot\zeta(\widehat Q\widehat P+\widehat P\widehat Q)/4\).}{Dans la paramétrisation elliptique, \(QP=I\sin2\theta\) et \(\dot\theta=-\omega\). Ainsi, \(H_{\rm tot}=I\omega+(I\dot\zeta/2)\sin2\theta\). La matrice quadratique a pour déterminant \(\omega^2-\dot\zeta^2/4\) : elle est définie positive précisément lorsque \(|\dot\zeta|<2\omega\), et semi-définie positive dans le cas d'égalité. La conservation de \(I\) et la dépendance temporelle de \(H_{\rm tot}\) expriment des propriétés différentes. À titre de contrôle algébrique, omettre le terme mixte donne \(\dot I=\dot\zeta(e^\zeta P^2-e^{-\zeta}Q^2)/2\). Dans la représentation quantique correspondante, le terme mixte symétrique est \(\dot\zeta(\widehat Q\widehat P+\widehat P\widehat Q)/4\).}

\subsection{\HMTLang{Compatibilité généalogique et transport conformément symplectique}{Compatibilité généalogique et transport conformément symplectique}}
\HMTLang{Dans la section engendrée,}{Sur la section engendrée,}
\[
 \zeta=\operatorname{artanh}\frac{2(\eta_{\rm ret}+\alpha)}{1000\alpha}.
\]
\HMTLang{Fixer \(\alpha\) et \(\eta_{\rm ret}\) fixe la forme. La version variable du théorème décrit le transport entre formes admises ; sa réalisation dynamique conserve la mise à jour qui détermine ces formes. Pour des sections positives \(\hbar_n\) dans une même coordinatisation dimensionnelle, soit \(\rho_n=\hbar_{n+1}/\hbar_n\). Alors}{Fixer \(\alpha\) et \(\eta_{\rm ret}\) fixe la forme. La version variable du théorème décrit le transport entre formes admissibles ; sa réalisation dynamique conserve la mise à jour qui détermine ces formes. Pour des sections positives \(\hbar_n\) dans la même coordinatisation dimensionnelle, soit \(\rho_n=\hbar_{n+1}/\hbar_n\). Alors}
\[
 \widetilde T_n=\sqrt{\rho_n}\,T_n,\quad
 \det\widetilde T_n=\rho_n,\quad
 \widetilde T_n^{\mathsf T}g_{\zeta_{n+1}}\widetilde T_n
 =\rho_ng_{\zeta_n},\quad
 \frac{I_{n+1}}{\hbar_{n+1}}=\frac{I_n}{\hbar_n}.
\]
\HMTLang{Ces égalités se déduisent en multipliant le transport précédent par \(\sqrt{\rho_n}\). La forme symplectique et l'action normalisée sont transportées conjointement. La limite différentiable a pour trace \(\dot\hbar/\hbar\) ; dans une forme symplectique fixe, un générateur hamiltonien linéaire a une trace nulle. La section variable exige donc de conserver la forme transportée ou les degrés de liberté qui échangent de l'action. Si la variation est un changement d'unités, on transporte d'abord les bases dimensionnelles.}{Ces égalités découlent de la multiplication du transport précédent par \(\sqrt{\rho_n}\). La forme symplectique et l'action normalisée sont transportées ensemble. La limite différentiable a pour trace \(\dot\hbar/\hbar\) ; pour une forme symplectique fixe, un générateur hamiltonien linéaire a une trace nulle. Une section variable exige donc de conserver la forme transportée ou les degrés de liberté échangeant de l'action. Si la variation est un changement d'unités, les bases dimensionnelles sont transportées en premier.}

\HMTLang{La distinction opérationnelle est précise. Sous un changement passif \(X=S_\zeta Y\), le lecteur satisfait lui aussi \(L_X(S_\zeta Y)=L_Y(Y)\) ; le terme mixte représente la connexion de coordonnées. Une déformation physique est spécifiée par l'évolution relative entre système et référence, obtenue à partir de leurs mises à jour et lue avec le même dispositif. Une attribution énergétique incorpore le champ, le registre et le contrôleur. Sont ainsi séparées la conservation démontrée de l'action quadratique, la réalisation hamiltonienne et les conditions d'une réponse physique vérifiable.}{La distinction opérationnelle est précise. Sous un changement passif \(X=S_\zeta Y\), le lecteur satisfait lui aussi \(L_X(S_\zeta Y)=L_Y(Y)\) ; le terme mixte représente la connexion de coordonnées. Une déformation physique est spécifiée par l'évolution relative entre système et référence, obtenue à partir de leurs mises à jour et lue avec le même dispositif. Une attribution énergétique comprend le champ, le registre et le contrôleur. Cela sépare la conservation démontrée de l'action quadratique, sa réalisation hamiltonienne et les conditions d'une réponse physique vérifiable.}
